\Needspace{11\baselineskip}
\SourceLeaf{277}
\begin{center}
\textit{Komposition der Klassen.}
\end{center}

\Article{249}
\textsc{Satz.} \textit{Gehören die Formen $f,f'$ jeweils zu denselben Ordnungen, Geschlechtern und Klassen wie $g,g'$, so gehört die aus $f$ und $f'$ zusammengesetzte Form zu derselben Klasse wie die aus $g$ und $g'$ zusammengesetzte.}

Aus diesem Satz, dessen Richtigkeit unmittelbar aus Artikel~239 folgt, ergibt sich die Bedeutung einer \textit{aus zwei gegebenen Klassen} oder auch \textit{aus mehreren zusammengesetzten Klasse}.

Wird irgendeine Klasse $K$ mit der Hauptklasse zusammengesetzt, so entsteht die Klasse $K$ selbst; die Hauptklasse kann also bei der Komposition mit anderen Klassen derselben Determinante weggelassen werden. Aus der Komposition zweier entgegengesetzter eigentlich primitiver Klassen entsteht stets die Hauptklasse derselben Determinante (siehe Artikel~243). Da jede ambige Klasse sich selbst entgegengesetzt ist, ergibt die Komposition irgendeiner eigentlich primitiven ambigen Klasse mit sich selbst die Hauptklasse derselben Determinante.

Der letzte Satz gilt auch umgekehrt: \textit{Ergibt die Komposition einer eigentlich primitiven Klasse $K$ mit sich selbst die Hauptklasse $H$ derselben Determinante, so muss $K$ eine ambige Klasse sein.} Ist nämlich $K'$ die zu $K$ entgegengesetzte Klasse, so ist die aus den drei Klassen $K,K,K'$ zusammengesetzte Klasse dieselbe wie die aus $H$ und $K'$ entstehende. Aus jenen entsteht $K$ (denn $K$ und $K'$ ergeben $H$, diese mit $K$ die Klasse $K$ selbst), aus diesen $K'$; somit fällt $K$ mit $K'$ zusammen und ist daher eine ambige Klasse.

Ferner ist der folgende Satz zu beachten: \textit{Sind die Klassen $K,L$ den Klassen $K',L'$ jeweils entgegengesetzt, so ist die aus $K$ und $L$ zusammengesetzte Klasse der aus $K'$ und $L'$ zusammengesetzten entgegengesetzt.} Die Formen $f,g,f',g'$ gehören jeweils zu den Klassen $K,L,K',L'$; die Form $F$ sei aus $f,g$ und $F'$ aus $f',g'$ zusammengesetzt. Da $f'$ der Form $f$ und $g'$ der Form $g$ uneigentlich äquivalent ist, während $F$ aus beiden $f,g$ direkt zusammengesetzt ist, ist $F$ auch aus $f',g'$ zusammengesetzt, aber aus beiden invers. Daher ist jede Form, die $F$ uneigentlich äquivalent ist, aus $f',g'$ direkt zusammengesetzt und somit $F'$ eigentlich äquivalent (Artikel~238,~239). Also sind $F,F'$ uneigentlich äquivalent, und die Klassen, zu denen sie gehören, sind entgegengesetzt.

Daraus folgt: Wird eine ambige Klasse $K$ mit einer ambigen Klasse $L$ zusammengesetzt. so entsteht stets eine ambige Klasse. Denn sie ist der Klasse entgegengesetzt, die aus den zu $K,L$ entgegengesetzten Klassen zusammengesetzt ist, also sich selbst, da diese Klassen sich selbst entgegengesetzt sind.

\SourceLeaf{278}
Schließlich bemerken wir: Sind zwei beliebige Klassen $K,L$ derselben Determinante gegeben, deren erste eigentlich primitiv ist, so lässt sich stets eine Klasse $M$ derselben Determinante finden, aus der zusammen mit $K$ die Klasse $L$ zusammengesetzt ist. Das erreicht man offenbar, indem man für $M$ die Klasse nimmt, die aus $L$ und der zu $K$ entgegengesetzten Klasse zusammengesetzt ist. Zugleich erkennt man sehr leicht, dass diese Klasse die einzige mit dieser Eigenschaft ist; verschiedene Klassen derselben Determinante ergeben also bei Komposition mit derselben eigentlich primitiven Klasse verschiedene Klassen.

Die Komposition der Klassen kann zweckmäßig mit dem Additionszeichen $+$ bezeichnet werden, ebenso wie die Gleichheit der Klassen mit dem Gleichheitszeichen. Mit diesen Zeichen lässt sich der eben gegebene Satz so darstellen: Ist $K'$ die zu $K$ entgegengesetzte Klasse, so ist $K+K'$ die Hauptklasse derselben Determinante, also $K+K'+L=L$; setzt man daher $K'+L=M$, so gilt $K+M=L$, wie verlangt. Gäbe es aber neben $M$ noch eine andere $M$ mit derselben Eigenschaft, also $K+M'=L$, so wäre $K+K'+M'=L+K'=M$ und damit $M'=M$. --- Werden mehrere identische Klassen zusammengesetzt, so kann man dies wie bei der Multiplikation durch Voranstellen ihrer Anzahl bezeichnen, sodass $2K$ dasselbe wie $K+K$, $3K$ dasselbe wie $K+K+K$ bedeutet usw. Dieselben Zeichen könnten auch auf Formen übertragen werden, sodass $(a,b,c)+(a',b',c')$ die aus $(a,b,c)$, $(a',b',c')$ zusammengesetzte Form bezeichnete. Um aber auch nur den Anschein einer Mehrdeutigkeit zu vermeiden, verzichten wir lieber auf diese Abkürzung, zumal wir einem solchen Zeichen $\surd M(a,b,c)$ bereits eine besondere Bedeutung gegeben haben. --- Wir sagen, die Klasse $2K$ entstehe durch \textit{Verdopplung} der Klasse $K$, die Klasse $3K$ durch \textit{Verdreifachung} usw.

\Needspace{8\baselineskip}
\Article{250}
Ist $D$ eine durch $mm$ teilbare Zahl, wobei wir $m$ als positiv voraussetzen, so gibt es eine Ordnung von Formen der Determinante $D$, die von der eigentlich primitiven Ordnung der Determinante $\frac{D}{mm}$ abgeleitet ist (oder \textit{zwei}, wenn $D$ negativ ist, nämlich eine positive und eine negative). Offenbar gehört die Form $(m,0,-\frac{D}{m})$ zu dieser Ordnung, und zwar zur positiven, und kann darin mit Recht als \textit{einfachste Form} betrachtet werden (ebenso ist $(-m,0,\frac{D}{m})$ die einfachste in der negativen Ordnung, wenn $D$ negativ ist). Gilt außerdem $\frac{D}{mm}\equiv1$ (modulo $4$), so gibt es auch eine Ordnung von Formen der Determinante $D$, die von der uneigentlich primitiven Ordnung der Determinante $\frac{D}{mm}$ abgeleitet ist; zu ihr gehört offenbar die Form $(2m,m,\frac{mm-D}{2m})$, die als einfachste in ihr gilt. (Ist $D$ negativ, so gibt es wiederum zwei Ordnungen, und in der negativen gilt die Form $(-2m,-m,\frac{D-mm}{2m})$ als einfachste.) So sind \textit{zum Beispiel}, wenn wir auch den Fall $m=1$ hier einbeziehen, in den vier Ordnungen
\SourceLeaf{279}
von Formen der Determinante $45$ die folgenden die einfachsten: $(1,0,-45)$, $(2,1,-22)$. $(3,0,-15)$, $(6,3,-6)$. Damit ergibt sich das folgende

\textsc{Problem.} \textit{Es sei irgendeine Form $F$ aus der Ordnung $O$ gegeben. Gesucht ist eine eigentlich primitive (positive) Form derselben Determinante, aus deren Komposition mit der einfachsten Form in $O$ die Form $F$ entsteht.}

\textit{Lösung.} Die Form $F=(ma,mb,mc)$ sei von der primitiven Form $f=(a,b,c)$ mit der Determinante $=d$ abgeleitet, und wir nehmen \textit{zuerst} an, dass $f$ eigentlich primitiv ist. Zunächst bemerken wir: Ist etwa $a$ nicht zu $2dm$ teilerfremd, so gibt es sicher andere, zu $(a,b,c)$ eigentlich äquivalente Formen, deren erste Koeffizienten diese Eigenschaft haben. Denn nach Artikel~228 gibt es zu $2dm$ teilerfremde Zahlen, die durch jene Form darstellbar sind. Es sei $a'=a\alpha\alpha+2b\alpha\gamma+c\gamma\gamma$ eine solche Zahl, und wir nehmen an, was erlaubt ist, dass $\alpha,\gamma$ teilerfremd sind. Wählt man dann $\beta,\delta$ so, dass $\alpha\delta-\beta\gamma=1$, so gehe $f$ durch die Substitution $\alpha,\beta,\gamma,\delta$ in die Form $(a',b',c')$ über; diese ist ihr eigentlich äquivalent und hat die verlangte Eigenschaft. Da auch $F$ und $(a'm,b'm,c'm)$ eigentlich äquivalent sind, erkennt man leicht, dass es genügt, den Fall zu betrachten, dass $a$ zu $2dm$ teilerfremd ist. Dann ist $(a,bm,cmm)$ eine eigentlich primitive Form (denn hätten $a,2bm,cmm$ einen gemeinsamen Teiler, so würde er auch $2dm=2bbm-2acm$ teilen) derselben Determinante wie $F$. Man bestätigt leicht, dass $F$ in das Produkt der Form $(m,0,-dm)$, die außer bei negativem $F$ die einfachste der Ordnung $O$ ist, mit $(a,bm,cmm)$ durch die Substitution $1,0,-b,-cm;\ 0,m,a,bm$ übergeht; daraus folgt nach dem Kriterium in Bemerkung~4 des Artikels~235, dass $F$ aus $(m,0,-dm)$ und $(a,bm,cmm)$ zusammengesetzt ist. Ist dagegen $F$ eine negative Form, so geht sie in das Produkt der einfachsten Form derselben Ordnung $(-m,0,dm)$ mit der positiven Form $(-a,bm,-cmm)$ durch die Substitution $1,0,b,-cm;\ 0,-m,-a,bm$ über und ist somit aus ihnen zusammengesetzt.

\textit{Zweitens}: Ist $f$ eine uneigentlich primitive Form, so darf man annehmen, dass $\frac{1}{2}a$ zu $2dm$ teilerfremd ist. Denn gilt diese Eigenschaft bei der Form $f$ noch nicht, so kann eine der Form $f$ eigentlich äquivalente Form mit dieser Eigenschaft gefunden werden. Daraus folgt leicht, dass $(\frac{1}{2}a,bm,2cmm)$ eine eigentlich primitive Form derselben Determinante wie $F$ ist. Ebenso leicht bestätigt man, dass $F$ in das Produkt der Formen
\[
(\pm2m,\pm m,\pm\tfrac{1}{2}(m-dm)),\quad
(\pm\tfrac{1}{2}a,bm,\pm2cmm)
\]
durch die Substitution
\SourceLeaf{280}
\[
1,0,\tfrac{1}{2}(1\mp b),-cm;\quad 0,\pm2m,\pm\tfrac{1}{2}a,(b\pm1)m
\]
übergeht. Dabei sind die unteren Zeichen zu nehmen, wenn $F$ eine negative Form ist, in den übrigen Fällen die oberen. Sie ist also aus diesen beiden Formen zusammengesetzt, deren erste die einfachste der Ordnung $O$, deren zweite eine eigentlich primitive (positive) Form ist.

\Needspace{8\baselineskip}
\Article{251}
\textsc{Problem.} \textit{Es seien zwei Formen $F,f$ derselben Determinante $D$ gegeben, die zu derselben Ordnung $O$ gehören. Gesucht ist eine eigentlich primitive Form der Determinante $D$, die mit $f$ zusammengesetzt $F$ ergibt.}

\textit{Lösung.} Es sei $\varphi$ die einfachste Form der Ordnung $O$; $\mathfrak{F}$, $\mathfrak{f}$ seien eigentlich primitive Formen der Determinante $D$, die mit $\varphi$ zusammengesetzt jeweils $F.\ f$ ergeben; schließlich sei $f'$ eine eigentlich primitive Form, die mit $\mathfrak{f}$ zusammengesetzt $\mathfrak{F}$ ergibt. Dann ist die Form $F$ aus den drei Formen $\varphi,\mathfrak{f},f'$ oder aus den beiden $f,f'$ zusammengesetzt. \textit{Q.~E.~I.}

Jede Klasse einer gegebenen Ordnung kann also als zusammengesetzt aus irgendeiner gegebenen Klasse derselben Ordnung und einer eigentlich primitiven Klasse derselben Determinante betrachtet werden.

\Needspace{11\baselineskip}
\begin{center}
\textit{Bei gegebener Determinante enthalten die einzelnen Geschlechter derselben Ordnung gleich viele Klassen.}
\end{center}

\Article{252}
\textsc{Satz.} \textit{Bei gegebener Determinante enthalten die einzelnen Geschlechter derselben Ordnung gleich viele Klassen.}

\textit{Beweis.} Die Geschlechter $G$ und $H$ sollen zu derselben Ordnung gehören; $G$ bestehe aus $n$ Klassen $K,K',K''\ldots K^{n-1}$, und $L$ sei irgendeine Klasse des Geschlechts $H$. Nach dem vorigen Artikel suche man eine eigentlich primitive Klasse $M$ derselben Determinante, aus deren Komposition mit $K$ die Klasse $L$ entsteht, und bezeichne die Klassen, die aus der Komposition von $M$ mit $K',K''\ldots K^{n-1}$ entstehen, jeweils mit $L',L''\ldots L^{n-1}$. Dann folgt aus der letzten Bemerkung des Artikels~249, dass alle Klassen $L,L',L''\ldots L^{n-1}$ verschieden sind; nach Artikel~248 gehören sie alle zu demselben Geschlecht, \textit{das heißt,}\ zum Geschlecht $H$. Schließlich erkennt man leicht, dass $H$ keine weiteren Klassen enthalten kann, da jede Klasse des Geschlechts $H$ als zusammengesetzt aus $M$ und einer anderen Klasse derselben Determinante betrachtet werden kann, die notwendig stets zum Geschlecht $G$ gehört. Somit enthält $H$ ebenso wie $G$ $n$ verschiedene Klassen. \textit{Q.~E.~D.}
